%% kapa: exam=1 round=22 year=2011 period=1 form=A subject=economics subject_ko=경제원론 seq=30 number=70
%% source: Q-Net 감정평가사 기출문제 (cst00309) artlSeq=5036790 file=extracted/1008/5036790/2011년 제22회 감정평가사 1차시험 기출문제/1교시.hwp
%% conversion: kordoc 4.24.51 -> tools/gen1.py (auto); review=proofread-sonnet; flags=none
\begin{question}{70}
총수요-총공급(AD-AS)모형에 관한 설명으로 옳지 않은 것은? (단, 총수요곡선은 우하향, 총공급곡선은 우상향)
\begin{choices}
  \item 독립투자 증가는 총수요곡선을 우측으로 이동시킨다.
  \item 정부지출 증가는 총수요곡선을 우측으로 이동시킨다.
  \item 조세 증가는 총수요곡선을 좌측으로 이동시킨다.
  \item 통화공급 증가는 총수요곡선을 좌측으로 이동시킨다.
  \item 기술 진보는 총공급곡선을 이동시킨다.
\end{choices}
\end{question}
\answer{4}
